\documentclass[12pt]{article}
\usepackage[T1]{fontenc}
\usepackage[utf8]{inputenc}
\usepackage{lmodern}
\usepackage{textcomp}
\usepackage{enumitem}
\usepackage{geometry}
\usepackage{graphicx}
\usepackage{amsmath, amssymb}
\geometry{margin=1in}
\begin{document}
\begin{center}
Biology - Undergraduate Level Assessment 25 \
Total Marks: 40 \
Answer all questions.
\end{center}

\section*{Multiple Choice (10 marks)}
\begin{enumerate}[label=\arabic*.]
\item Using conditioning, how can one determine the sensitivity of a dog's hearing to sounds of different frequencies?
\begin{enumerate}[label=(\alph*)]
    \item By using a controlled group of dogs not exposed to the sounds
    \item By using sensory deprivation to heighten hearing sensitivity
    \item By observing the dog's natural behavior in different sound environments
    \item By using negative reinforcement only
    \item By using operant conditioning
\end{enumerate}
\item How can one demonstrate that amino acids are used tobuild proteinsin cells, while the proteins already in the cell are usedin catabolism?
\begin{enumerate}[label=(\alph*)]
    \item Enzyme-linked immunosorbent assay (ELISA)
    \item Polymerase chain reaction (PCR)
    \item X-ray crystallography
    \item Mass spectrometry
    \item Autoradiography
\end{enumerate}
\item What prevents the stomach from being digested by itsown secretions?
\begin{enumerate}[label=(\alph*)]
    \item The thick muscular wall of the stomach
    \item Antibacterial properties of stomach acid
    \item Gastric juice
    \item High pH levels in the stomach
    \item Presence of beneficial bacteria
\end{enumerate}
\item Why would you not expect conjugation among a group of paramecia that had descended from a single individual through repeated fission?
\begin{enumerate}[label=(\alph*)]
    \item Conjugation is prevented by the presence of a dominant allele
    \item Conjugation results in lesser genetic diversity
    \item Conjugation can occur among identical organisms
    \item All individuals are identical genetically, and conjugation could not occur.
    \item Conjugation does not require different mating types
\end{enumerate}
\item If the activity of an enzyme is constant over a broad range of pH values, it is likely that
\begin{enumerate}[label=(\alph*)]
    \item the enzyme is inactive at all pH values
    \item no ionizing groups on the enzyme or substrate participate in the reaction
    \item only ionizing groups on the substrate participate in the reaction
    \item ionizing groups on both the enzyme and substrate participate in the reaction
    \item only non-ionizing groups on the substrate participate in the reaction
\end{enumerate}
\end{enumerate}

\section*{True / False (10 marks)}
\textit{Answer True or False}
\begin{enumerate}[label=\arabic*.]
\setcounter{enumi}{5}
\item True or False: Aposomatic coloration is a strategy where organisms use camouflage to avoid detection by predators.
\item True or False: The ability of the brain to detect differences in stimulus intensity is influenced by the number of synapses crossed by sensory signals.
\item True or False: The synthesis of an RNA/DNA hybrid from a single-stranded RNA template requires a DNA or RNA primer and reverse transcriptase.
\item True or False: The process of molting in crustaceans is controlled by hormonal influences rather than just water temperature alone.
\item True or False: Hemoglobin's affinity for O$_{2}$ increases as blood pH decreases, enhancing oxygen delivery to tissues.
\end{enumerate}

\section*{Long Form Response (20 marks)}
\textit{Answer each question with a clear and complete explanation.}
\begin{enumerate}[label=\arabic*.]
\setcounter{enumi}{10}
\item Discuss the role of water availability in shaping the transition from forest ecosystems to grasslands as one moves from east to west across the U.S.
\item Discuss the concepts of kinesis and taxis in animal behavior, highlighting their differences and providing examples of each in natural settings.
\end{enumerate}

\vfill
\noindent\textit{End of Paper}
\end{document}